%%%%%%%%%%%%%%%%%%%%%%%%%%%%%%%%%%%%%%%%%%%%%%%%
% COPYRIGHT: (C) 2012-2026 FAU FabLab and others
%  CC-BY-SA 3.0
%%%%%%%%%%%%%%%%%%%%%%%%%%%%%%%%%%%%%%%%%%%%%%%%


\newcommand{\basedir}{./fablab-document/}
\documentclass{\basedir/fablab-document}
\usepackage{tabularx} % Tabellen mit bestimmtem Breitenverhältnis der Spalten
\usepackage{ifthen}
\usepackage{xspace}
\def\tabularnewcol{&\xspace} % hässlicher Workaround von http://tex.stackexchange.com/questions/7590/how-to-programmatically-make-tabular-rows-using-whiledo

\title{Einweisungsliste Shaper Origin}
\fancyfoot[L]{Einweisungsliste Nr. \underline{\hspace{3em}}}
% Versionsnummer: version.tex wird von der GitHub Action erzeugt,
% lokal steht stattdessen "Entwurf" mit Datum in der Fußzeile.
\InputIfFileExists{version}{}{\newcommand{\Version}{Entwurf \today}}
\fancyfoot[R]{Version \Version}

\begin{document}
\vspace*{-2em}

\textbf{für die handgeführte CNC-Fräse Shaper Origin}

Ich bestätige mit meiner Unterschrift verbindlich, dass

\begin{itemize}
\item ich die Regeln und Hinweise, die Betriebsanweisung sowie die Kurzanleitung des Herstellers
gelesen und verstanden habe
\item ich ein Werkstück unter Anleitung erfolgreich gefräst habe, mit Befestigen des Werkstücks,
Kleben und Scannen des ShaperTape, Fräserwechsel und Z-Touch
\item ich bei weitergehenden Fragen die Anleitung des Herstellers lesen werde
\item die Möglichkeit für Fragen und Diskussion gegeben war
\item die Möglichkeit, eine Kopie der Anleitung zu erhalten, gegeben war
\end{itemize}



% einfach kopiert von Einweisungsliste Lasercutter

\newcounter{i}
\setcounter{i}{1}

\newcommand{\leerezeile}{\hspace{2em} \tabularnewcol \hspace{3em} \tabularnewcol \hspace{2.5em} \tabularnewcol \hspace{2.5em} \tabularnewcol \vbox{\vspace{2em}} \tabularnewcol \tabularnewcol \tabularnewcol \tabularnewline \hline}

\begin{tabularx}{\textwidth}{|l|l|l|l|X|X|X|X|}
  \hline
  \textbf{Nr.} & \textbf{Datum} & \textbf{von} & \textbf{bis} & \textbf{Name} & \textbf{Unterschrift} & \textbf{Einweisender} & \textbf{Unterschrift} \\ \hline
  \whiledo{\value{i}<12}%
  {%
    \stepcounter{i} \leerezeile
  }%
  \leerezeile % doofer Workaround, eigentlich sollte das auch in der Forschleife gehen! Ohne dies wird die Spaltenbegrenzung von Spalte 1 zu weit gezeichnet.
\end{tabularx}

\end{document}
